% Draft proof unit for YM-000026. Parent supplies theorem environments/macros.
\section{Finite-volume stress shell nodes}
\label{sec:ex-lattice-actual-stress-shell-nodes}

This section records the actual finite-volume producer used by the stress
annulus argument. It follows the selected family and its retained history.
The scale choice is made before the reference opening, insertion labels, or
finite list of tests. The quantitative conclusion at a fixed collar site is
then obtained from a two-channel coefficient construction and the finite
quota-to-norm comparison.

Write $a_i$ for the physical spacing, $\lambda_i=\log(a_i^{-1})$,
$B_i$ for the support radius, and $\widehat\beta_i$ for the effective
coupling. Put
\[
 \begin{aligned}
 r_m&=128 B_{m+1}^{2}a_{m+1}, & c_m&=\exp(\lambda_m/32),\\
 X_m&=\exp\sqrt{\widehat\beta_{m+1}}, & n_m&=\text{the number of columns at level }m.
 \end{aligned}
\]
For a fixed finite-volume law, let $\mathcal O_{N,q}(L,x)$ be the readout
of label $L$ at site $x$. For an opened reference level $p_0$, let
$\mathcal C_{N,q}(p_0)$ be the finite cell index set of the centered chamber
tiling used by the root test, and write $F_v$ for the prescribed translated
copy of $F$ at $v\in\mathcal C_{N,q}(p_0)$. The copy includes complex
conjugation exactly when the cell has odd reflection parity. Set
\begin{equation}\label{eq:ex-lattice-shell-root-functional}
 \mathcal R_{N,q,p_0}(F)=
 \left|\mathbb E_{\beta,N,q}
      \prod_{v\in\mathcal C_{N,q}(p_0)}F_v\right|^{1/|\mathcal C_{N,q}(p_0)|}.
\end{equation}
This is the root-test functional of that opening, not a product over every
fine-lattice site. The grade-$D$ quota at birth level $m$ is denoted
$\mathfrak q_{D,m,N,q}(z)$. It is a coefficient size in the fixed reference
frame; it is not the norm in the continuum field space. The opening $p_0$
selects the root-test chamber and its stake bound; the quota index is $m$.

For grades $d,D,h\in\mathbb N$ satisfying
\begin{equation}\label{eq:ex-lattice-shell-grade}
 D+h=d+4,
\end{equation}
and labels $B$ of grade $d$ and $T$ of grade four, define the finite
fragment
\begin{equation}\label{eq:ex-lattice-shell-fragment}
 \mathcal F_{m,x}^{N,q}(z)
 =r_m^h\mathcal O_{N,q}(T,x_{m,N,q})\mathcal O_{N,q}(B,0)
       -\mathcal O_{N,q}(z,0),
 \qquad z\text{ of grade }D.
\end{equation}
Here $x_{m,N,q}$ is the image of the fixed cover anchor $x$ under the
volume-period placement at level $j_m$. The factor is $r_m^h$, where
$r_m=128B_{m+1}^2a_{m+1}$; it is not $a_{m+1}^h$. Identity
\eqref{eq:ex-lattice-shell-grade} makes both terms have the same scaling
power (the labels themselves lie in filtered grade spaces).
Write $\operatorname{cs}_q(v)$ for the coefficient size of a grade-$q$
label in the selected frame.

\subsection{The scale schedule and its numerical gates}

Fix a simple compact gauge group, the selected Wilson family, and its
same-history records. The scale estimates used by the source give the
following schedule for each fixed triple $(d,D,h)$ satisfying
\eqref{eq:ex-lattice-shell-grade}:
\begin{equation}\label{eq:ex-lattice-shell-schedule}
 t=1+\frac{1}{64(d+12)},\qquad
 j_m\longrightarrow\infty,\qquad
 \frac{\lambda_{j_m}}{\lambda_m}\longrightarrow t.
\end{equation}
The sequence is chosen independently of the reference opening and all
labels. The strict rate condition is
\[
 (d+12)(1-t)=-\frac1{64}>-\frac1{32}.
\]
The separate lightness estimate gives, for every fixed opening $p_0$,
eventually in $m$,
\begin{equation}\label{eq:ex-lattice-shell-lightness}
 32r_m\le a_{p_0},\qquad
 a_{m+1}\bigl(B_{m+1}^{2}+1+2\cdot1024B_{m+1}^{2}\bigr)
       \le a_{p_0}.
\end{equation}
The first inequality supplies the collar margin. The second supplies the
number of centered chamber cells needed by the finite stencil.

Here are the channel prices in ordinary notation. Write $\delta_i$ for the
renal allowance, $\phi_i$ for the near allowance, and $s_i$ for the support
exponent. Set
\[
 q_i=(2\cdot1024B_i+1)^4,\qquad
 \operatorname{krill}_{j,i}(u)=\frac{u}{16}\left(\frac{a_j}{a_i}\right)^4.
\]
The quantity $q_i$ is the exponential-count parameter; it is distinct from
the geometric radius $1024B_i^2$ used in the finite stencil. Let $g_q$ be
the maximal engineering degree at grade $q$. The paid grade-$q$ weight for
a label-size bound $v>0$ is
\[
 \mathsf w_q(v,j)=\frac{\widehat\beta_j^{-(g_q+4)}}
  {v\,C_j\prod_{\ell<j}C_{\ell+1}},
\]
where $C_i$ is the column count. Set
\[
 A_-=a_{j_m}^{d}\mathsf w_d(a,j_m),\qquad
 A_+=r_m^{-h}a_{j_m}^{4}\mathsf w_4(b,j_m).
\]
For channel weights $A_-,A_+$, define
\begin{equation}\label{eq:ex-lattice-shell-picket}
 \rho_{m,\pm}=\frac{\delta_{m+1}^{2}
   \operatorname{krill}_{j_m,m+2}
      (\delta_{j_m}\phi_{j_m}A_\pm)}
      {\exp((s_{m+1}+1/2)q_{m+1})},\qquad
 W_m=\rho_{m,-}\rho_{m,+}.
\end{equation}
This follows from the scalar exponential weight in each channel at
$i=m+1$: both outer and inner renal factors are $\delta_{m+1}$, and the
input also contains the near allowance $\phi_{j_m}$. The channel weights $A_-$
and $A_+$ are the paid prices for the grade-$d$
and grade-four insertions; the latter includes $r_m^{-h}$, cancelled by
$r_m^h$ in the fragment. Thus $W_m$ is a positive product fixed by the
source formulas, not a free normalization parameter.

For a positive function $w$ on channel names, define its marched price
from level $k$ by
\[
 (\mathsf M_{k,m}w)(i)=
 \begin{cases}
 w(i),&m<k,\\
 w(i)\displaystyle\prod_{s=k}^{m-1}\phi_s,&m\ge k,
 \end{cases}
\]
where $0<\phi_s\le1$ is the selected transition allowance. These formulas
are the two positive channel prices; in particular $W_m$ is not a free
normalization parameter.
The selected price estimate supplies $W_m>0$ and
\begin{equation}\label{eq:ex-lattice-shell-march}
 \mathsf M_{k,m}(c_m)\le \delta_m\frac{W_m}{16X_m}.
\end{equation}
Set
\begin{equation}\label{eq:ex-lattice-shell-coefficient-budget}
 e_m=32(1+n_m)a_{m+1}^{D}\frac{X_m}{W_m}.
\end{equation}
The second price estimate, with
$\|a_k^{-D}\|=a_k^{-D}$ and
$g_{m,k}=\exp((\lambda_m-\lambda_k)/8)$, is
\begin{equation}\label{eq:ex-lattice-shell-correction-price}
 32(1+n_m)c_m a_k^{-D}\frac{e_m}{g_{m,k}}
       \le \frac12\widehat\beta_k^{-\nu_D},
 \qquad \nu_D=g_D+4.
\end{equation}
Here $g_D$ is the maximal engineering degree of grade-$D$ labels. The
normalization on the right is the grade-$D$ payment. For fixed
$(p_0,k)$ it is independent of $m$.
To expose the rate calculation, for a positive sequence $f_m$ write
\[
 \operatorname{rate}(f)=\lim_{m\to\infty}\frac{\log f_m}{\lambda_m}
\]
when the limit exists. The rate rules for the two channel weights give
\[
 M=(d+12)(1-t)=-\frac1{64},\quad
 \operatorname{rate}(W_m)=-D+M.
\]
Explicitly, the unscaled grade-$d$ weight has rate $-dt$ and the unscaled
grade-four weight has rate $-4t+h$, since it is divided by $r_m^h$ and
$\operatorname{rate}(r_m)=-1$. Each channel transport contributes
$4(1-t)$ from its fourth-power block ratio. Thus
\[
 \operatorname{rate}(\rho_{m,-})=-dt+4(1-t),\qquad
 \operatorname{rate}(\rho_{m,+})=-4t+h+4(1-t),
\]
and their sum is $h+8-(d+12)t=-D+(d+12)(1-t)$, using
$D+h=d+4$. The renal and near allowances have rate zero. The explicit
source bound $(s_i+1/2)q_i\leq\sqrt{\widehat\beta_{i+1}}$ shows that the
exponential denominator has rate zero as well. In the displayed paid width,
$\log\widehat\beta_j=o(\lambda_j)$ and
$\sum_{\ell<j}\log C_{\ell+1}=o(\lambda_j)$, so its coupling and column
factors have rate zero. Also $\operatorname{rate}(c_m)=1/32$,
$\operatorname{rate}(a_{m+1}^D)=-D$, and
$\operatorname{rate}(X_m)=0$, because
$X_m=\exp(\sqrt{\widehat\beta_{m+1}})$ is subexponential in $\lambda_m$.
The finite column census is likewise subexponential. The exact definition
of $e_m$ therefore gives $\operatorname{rate}(e_m)=-M=1/64$.
In particular $c_m^{-1}\to0$, since $\lambda_m\to\infty$.
For fixed $k$, $g_{m,k}=\exp((\lambda_m-\lambda_k)/8)$ has rate $1/8$,
while $a_k^{-D}$ and $\widehat\beta_k^{-\nu_D}$ are fixed positive
constants. Thus the correction-price left side has rate
\[
 \frac1{32}-M-\frac18=-\frac5{64}<0;
\]
it tends to zero and is eventually below the displayed grade-$D$
payment. For the marched estimate, the one-step growth bound and
$\operatorname{fly}(D,m)\leq1$ give
\[
 \mathsf M_{k,m}(c_m)\leq
 \frac{c_m}{\upsilon_k}\left(\frac{a_m}{a_k}\right)^{D+1},
 \qquad \upsilon_i=\operatorname{fly}(D,i),\quad k\geq\operatorname{ref}(D).
\]
For fixed $k$, the right side has rate $1/32-(D+1)$, while
$\delta_mW_m/(16X_m)$ has rate $-D+M$. Their rate difference is
$1/32-1-M=-61/64<0$, so the former is eventually smaller. This uses
the same-history growth estimate, the rate-zero allowance, and the paid
channel-rate identity above. The separate strict condition $M>-1/32$
is a sufficient source condition for the correction price to be eventual.
Hence both price
inequalities hold for each fixed opening.

The other price gate used in the finite coefficient construction is the exact budget
identity
\begin{equation}\label{eq:ex-lattice-shell-ledger-budget}
 4(1+n_m)\,2\left(a_{m+1}^{D}X_m\frac{4}{W_m}\right)=e_m.
\end{equation}
Thus the available coefficient budget is precisely the finite-census
price of the principal term, with its two channel factors; no extra
polynomial is hidden in this equality.

\subsection{The coefficient and density ledger}

Set
\begin{equation}\label{eq:ex-lattice-shell-one-step-scales}
 a_m^*=\frac4{W_m},\qquad I_m=\frac{W_m}{16X_m},\qquad
 p'_m=a_{m+1}^{D}(X_m a_m^*),\qquad e_m=8(1+n_m)p'_m.
\end{equation}
Then $e_m=32(1+n_m)a_{m+1}^D X_m/W_m$, as above. This reparameterization
exhibits the numerical inputs to the coefficient construction. Put
$L_m=a_m/a_{m+1}$ and $A_m=L_m^D$, the block dilation of a grade-$D$
coefficient. The block
side $L_m\geq2$ and grade $D\geq0$ give $\|A_m^{-1}\|=L_m^{-D}\leq1$.
The stretch across this step satisfies
\[
 a_m^D L_m^{-D}=a_{m+1}^D,
\]
so the coefficient bound $X_m a_m^*$ delivered by the finite-coupon
construction becomes exactly $p'_m=a_{m+1}^D(X_m a_m^*)$ after this
grade-$D$ rescaling.
Together with $X_m\geq1$,
\[
 2I_m(2a_m^*)=X_m^{-1}\leq1,\qquad
 I_m\|A_m^{-1}\|(X_ma_m^*)=\frac{\|A_m^{-1}\|}{4}\leq1.
\]
The marched price says $J_m:=\mathsf M_{k,m}(c_m)\leq\delta_m I_m$.
At each name the selected menu has the corresponding regular price
$J_m$ and allowance $\delta_m I_m$; multiplying the finite factors and
using $1+v\leq e^v$ for $v\geq0$ yields the finite transfer bound
with its exponential penalty. This is the
finite-product estimate used by the transition, not a claim that normalized
root expectations multiply.

The fine response is first built for each of the two
insertions from the same terminal family, door, and funded path. Its quiet
coefficient is the background activity and its singleton coefficient is
the paid target plus the hook correction; signed covariance transports
these identities from the distinguished name to the other names. Backward
transport uses
$\operatorname{krill}_{j,i}(u)=u(a_j/a_i)^4/16$. If
$L_i=a_i/a_{i+1}\geq2$, then
\[
 \operatorname{krill}_{j,i+1}(u)=L_i^4\operatorname{krill}_{j,i}(u),
 \qquad 6L_i\operatorname{krill}_{j,i}(u)
       \leq\operatorname{krill}_{j,i+1}(u),
\]
because $L_i^3\geq8$. Iteration moves the response to level $m+1$;
injective migration preserves occupancy and the block map expands
separation by at least $L_i$. The fine response uses occupancy
$o_{m+1}=r(m+1)$.

The two channels refer to the two inserted labels. Their graft is formed at the same base name, and source affine
covariance fixes its quiet and singleton coefficients. The higher
many-name terms are pruned with the finite-mask receipt. With fine
occupancy $o_{m+1}$ its mask cost is
\[
 \sigma_m=2(1+o_{m+1})\eta_{m+1}
       \bigl(o_{m+1}e^{s_{m+1}+1/2}+1\bigr)
       +(1+o_{m+1})\chi_{m+1}.
\]
The selected finite-mask receipt gives
$(2B_{m+1}+1)^4\sigma_m\leq1/2$ and
$e^{\sigma_m}-1\leq1$. The radius $1024B_{m+1}^2$ in the stencil and
the exponent count $q_{m+1}$ in the channel price are distinct quantities.

The finite-coupon step and transition then yield a coefficient $Q_m$ at
level $m$. Its principal coefficient is bounded by $p'_m$. The census
cost is exactly
\[
 4(1+n_m)\,2p'_m=e_m,
\]
where $n_m$ is the column count; the factor $4$ is the source's two-channel
and sign count. The correction has the purchased one-step remainder
bound. No response coefficient is declared exact before that paid
correction is included.

Here is the exact density identity retained by the construction. The
paid representation supplies regular coefficients $H_S(U)$ indexed by
supports $S$ and forced coefficients $K_P(U)$ indexed by boxes $P$. Each
regular coefficient is supported on its assigned $S$, and each forced
coefficient on its assigned box support. Write $E_m(U)$ for the action
polynomial, $\Phi_i(U)$ for singleton field coefficients, $\bar A_P(U)$
for the background activity, and $\omega_{i,P}(U)$ for singleton
activities. The operator $\operatorname{send}_i(z)$ is the polynomial
monomial with coefficient $z$ in insertion slot $i$; $F(U)$ is the target
flag set. Then the pushed density is
\begin{align}\label{eq:ex-lattice-shell-served-density}
 \mathscr D_m(U)={}&e^{E_m(U)}\exp_{\mathrm{nil}}\!\left(
       \sum_i\operatorname{send}_i(\Phi_i(U))+\sum_S H_S(U)\right)\\
 &\quad\times
 \sum_{\substack{\mathcal P\text{ a family of boxes}:\\
        P\cap P'=\varnothing\ (P\ne P'),\quad
        \bigcup_{P\in\mathcal P}Q_P=F(U)}}
       \prod_{P\in\mathcal P}
       \left(\bar A_P(U)+\sum_i\operatorname{send}_i(\omega_{i,P}(U))
                    +K_P(U)\right).
\end{align}
Each box is $(P,Q_P)$ with $Q_P\subseteq P$; compatibility means
disjointness of full supports, while the flags cover $F(U)$. Here
$\exp_{\mathrm{nil}}(Z)=\sum_{r=0}^{|I|}Z^r/r!$ for the finite set $I$ of
insertion names. The degree conditions make every $H_S$ and $K_P$ vanish
in degrees zero and one; their regular-support and forced-box coefficient
costs have separate summable bounds.
Thus the higher-degree exponential terms, forced terms, and compatible
box products remain explicitly present and paid. In particular, the
pushed density is not identified with a single leading coefficient.
The canonical local-transfer proof pays these two arrays separately and
proves the same-law density equality used here.

The served transfer uses coarse occupancy $o_m=r(m)$; the fine response
and graft use $o_{m+1}=r(m+1)$. Set
$d_m=\widehat\beta_{m+1}^{-2}$,
$u_m=\widehat\beta_{m+1}^{-1}$, and $y_m$ equal to the flagged response
price. For the explicit rooted-site cost, put
$\mathsf S(B)=e(2(16B+9))^4$. The transfer duty is
\[
 0<d_m\leq1,\qquad
 T_m=o_m\mathsf S(B_{m+1})d_m+4d_m,\qquad
 T_m+e^{T_m}T_m\leq u_m,
\]
together with
\[
 (2B_m+1)^4\bigl(y_m+o_m\mathsf S(B_m)d_m\bigr)
       \leq\tfrac12,\qquad
 e^{y_m+o_m\mathsf S(B_m)d_m}-1\leq1.
\]
Here $y_m=2\exp(-\mu_m\Delta_{G,L_m}(\widehat\beta_m)^2/16)$,
where $\mu_m,\Delta_{G,L_m}$ are the selected shape strength and fine
monitor. These are the exact coarse transfer inequalities; the fine
backward transport and graft instead use $o_{m+1}$. The endpoint support
exponents are at least $100$. The source tariff and same-history service
proofs in Sections~\ref{sec:finite-menu-history-handoff} and
\ref{sec:local-transfer} supply these paid conditions. This keeps the
regular and forced coefficient charges separate and does not assume
equality of the interacting and reference actions.

\subsection{From the ledger to one collar node}

Fix an opened reference level $p_0$ and a level $k$ with that opening,
$k$ beyond the grade-$D$ threshold, and labels $B,T$. The quota estimate
at birth level $m$ gives a finite constant $Q_m^{\mathrm{quota}}$
independent of volume and period. It may depend on $m$; no uniform
birth-level bound is asserted:
\begin{equation}\label{eq:ex-lattice-shell-quota}
 \|z\|_D\le Q_m^{\mathrm{quota}}\mathfrak q_{D,m,N,q}(z).
\end{equation}
Apply the ledger and the same-history coefficient-to-label transfer at a
fixed collar anchor. The coefficient budget
\eqref{eq:ex-lattice-shell-ledger-budget} and the finite census give
$\mathfrak q_{D,m,N,q}(z)\le2e_m$, hence
\[
 \|z\|_D\le2Q_m^{\mathrm{quota}}e_m.
\]
The transferred label satisfies the exact product-minus-correction
identity \eqref{eq:ex-lattice-shell-fragment}. Its chamber condition follows
from the two input chamber memberships and closure under products and
subtraction of a paid correction. Write $\Theta_k$ for the selected
treasury at level $k$. The response estimate together with
\eqref{eq:ex-lattice-shell-correction-price} gives
\begin{equation}\label{eq:ex-lattice-shell-root}
 \mathcal R_{N,q,p_0}(\mathcal F_{m,x}^{N,q}(z))
       \le 2S(p_0,k)c_m^{-1},
 \qquad
 S(p_0,k)=\exp\!\left((a_{p_0}/a_k)^4\Theta_k+5\right).
\end{equation}
To see the factor $2$, write the pre-correction function as $f$ and the
finite-coupon correction as $r'$. The exact decomposition is
\[
 f-\mathcal O_{N,q}(B,0)
   =(f+r')+(-r'-\mathcal O_{N,q}(B,0)).
\]
Put $B_m(p_0,k)=S(p_0,k)c_m^{-1}$. The two root-test bounds are respectively
at most $B_m(p_0,k)$ and $B_m(p_0,k)$ (the second is first obtained with the
stronger bound $B_m(p_0,k)/2$). The root-test triangle inequality applies because both
terms are in the centered chamber. Thus the factor $2$ comes from this
two-term triangle inequality; it is not obtained by multiplying normalized
root functionals or by combining regular and forced charges. The estimate uses
the served coefficient identity, the throw coupon, the fixed opening, and
the two inequalities \eqref{eq:ex-lattice-shell-march} and
\eqref{eq:ex-lattice-shell-correction-price}.

The placement and room thresholds are pointwise in the cover anchor. The
anchor is assumed to lie in the stated collar. For fixed $x$, the period
eventually places it at its intended grid site, so the radial bounds carry
over to the finite-volume location. The first inequality in
\eqref{eq:ex-lattice-shell-lightness} is the shell-smallness condition
used with those bounds to keep the remote insertion in the centered
chamber; the second gives enough centered chamber cells for the finite
response stencil. At each $(N,q)$ on that anchor's
late set, the construction just described gives its correction. Outside
that late set choose the zero label. The numerical bounds do not depend on
$x$, although the late set may.

\begin{proposition}[Selected-history shell-node output]
\label{prop:ex-lattice-actual-shell-node}
Take the selected public history of Theorem~\ref{thm:ex-selected-run-service}.
Use its adjacent-scale paid menu, tariff, and served-transfer receipts from
Sections~\ref{sec:ex-lattice-rg-multisource},
\ref{subsec:ex-lattice-volant-producer},
\ref{sec:finite-menu-history-handoff}, and \ref{sec:local-transfer}, together
with the scale and price gates displayed above and the quota estimate at
each birth level. For every fixed grade triple
$D+h=d+4$, there is a sequence $j_m$ independent of $p_0,k,B,T$ and of
the finite test family, with $\lambda_{j_m}/\lambda_m\to t$ as in
\eqref{eq:ex-lattice-shell-schedule}. For each fixed opened $(p_0,k)$ and
grades $d,D,h$, choose labels $B,T$ and positive coefficient bounds
$a,b$ with $\operatorname{cs}_d(B)\le a$ and $\operatorname{cs}_4(T)\le b$,
along with the finite extraction test family. Assume the selected stock
opens $p_0\leq k\leq m$, with $\operatorname{ref}(D)\leq k$,
$m+1<j_m<N$, and $\operatorname{ref}(d),\operatorname{ref}(4)\leq j_m$.
Assume also the same-history
chamber and room receipts: $\mathcal O_{N,q}(B,0)$,
$\mathcal O_{N,q}(T,\xi_m(x))$, and $\mathcal O_{N,q}(z,0)$ for every
grade-$D$ label $z$ lie in the centered chamber, and
\[
 32r_m\leq a_{p_0},\qquad
 a_{m+1}(B_{m+1}^2+1+2\cdot1024B_{m+1}^2)\leq a_{p_0}.
\]
There is
$m_0$ such that for $m\ge m_0$ and each fixed cover anchor $x$ whose
intended level-$j_m$ physical location $\xi_m(x)$ obeys
$r_m/16\le\|\xi_m(x)\|_2\le4r_m$, the finite-volume correction has
\[
 \|z_x(N,q)\|_D\le2Q_m^{\mathrm{quota}}e_m,
 \qquad \mathcal R_{N,q,p_0}(\mathcal F_{m,x}^{N,q}(z_x(N,q)))
       \le2S(p_0,k)c_m^{-1}
\]
for filter-almost every pair along the selected volume and period filters. The
late event may depend on $x$; the displayed numeric bounds do not.
\end{proposition}

\begin{proof}
Choose $t$ and $j_m$ by the rate schedule. The strict inequality
$(d+12)(1-t)>-1/32$ makes both rate-based price gates eventual. For fixed
$p_0$, the two lightness inequalities are eventual as well. Thus for each
fixed data set choose $m$ beyond the common numeric thresholds, the grade
opening, and the finite-stencil room threshold. At a fixed cover anchor,
choose the late volume-period event on which its image is the prescribed
grid site and the opening is present. The local ledger gives the correction
and its chamber/root bounds; the fixed quota estimate converts its quota
to the continuum coefficient norm. This produces exactly the pointwise
late-set quantifier stated above. It does not claim one late event
simultaneously for the growing collar.
\end{proof}

\subsection{Scope of the shell construction}

The menu and transfer receipts invoked above are the same-history
constructions of Sections~\ref{sec:ex-lattice-rg-multisource},
\ref{subsec:ex-lattice-volant-producer},
\ref{sec:finite-menu-history-handoff}, and
\ref{sec:local-transfer}; the selected public service and its endpoint
support exponents are supplied by Theorem~\ref{thm:ex-selected-run-service}.
The present argument supplies the remaining shell-specific composition:
the two transported channel responses, finite mask and coupon steps,
coefficient ledger, placement, and collar norm/root bounds. The collar's
late event is pointwise in its fixed anchor, as required by the source.
